% !TEX root = ../main.tex

\section{Vectors and embeddings}

\begin{frame}{Scalar, vector, matrix, tensor}
  \small
  These names describe how numbers are arranged into axes:

  \vspace{0.7em}
  \begin{tabular}{@{}p{2.0cm}p{4.7cm}p{3.9cm}@{}}
    \toprule
    \textbf{Name} & \textbf{Example} & \textbf{Rank and shape} \\
    \midrule
    Scalar & \(s=3.2\) & rank 0, shape \(()\) \\[0.65em]
    Vector & \(\mathbf{x}=\begin{bmatrix}0.2&-0.7&0.4\end{bmatrix}\)
      & rank 1, shape \((3)\) \\[0.65em]
    Matrix & \(X=\begin{bmatrix}0.2&-0.7&0.4\\0.1&0.8&-0.3\end{bmatrix}\)
      & rank 2, shape \((2,3)\) \\[0.65em]
    Tensor & \(T=\left[\begin{matrix}X_1\\X_2\\\vdots\end{matrix}\right]\)
      & rank 3, shape \((b,n,d)\) \\
    \bottomrule
  \end{tabular}

  \vfill
  \begin{block}{Terminology}
    A tensor is the general object; a scalar, vector, and matrix are rank-0,
    rank-1, and rank-2 tensors. Shape gives the size along every axis.
  \end{block}
\end{frame}

\begin{frame}{Embedding lookup: from an ID to a vector}
  \small
  Let \(|V|\) be the number of tokens. The embedding dimension \(d\) is a model
  hyperparameter: the chosen width of every token vector. The model stores the
  learnable \emph{embedding matrix} \(W_E\):
  \[
    W_E\in\mathbb{R}^{d\times |V|}.
  \]

  \begin{center}
    \begin{tikzpicture}[>=Latex]
      \node[draw=tokenblueborder, fill=tokenblue, rounded corners,
        minimum width=1.6cm, minimum height=0.8cm] (id) {ID \(i=2\)};

      \node[right=2.2cm of id] (matrix) {\(
        W_E=\begin{bmatrix}
          0.1 & \phantom{-}0.7 & \colorbox{tokengreen}{\(-0.3\)} & 0.0\\
          -0.4 & \phantom{-}0.2 &
          \colorbox{tokengreen}{\(\phantom{-}0.8\)} & -0.2\\
          0.6 & -0.1 &
          \colorbox{tokengreen}{\(\phantom{-}0.5\)} & \phantom{-}0.9
        \end{bmatrix}\)};

      \node[draw=tokengreenborder, fill=tokengreen, rounded corners,
        right=1.2cm of matrix, align=center] (vector)
        {\(\mathbf{e}_i=(W_E)_{:,i}\)\\[0.2em]
         \(=[-0.3,0.8,0.5]^{\mathsf T}\)};

      \draw[->, thick] (id) -- node[above]{select column} (matrix);
      \draw[->, thick] (matrix) -- (vector);
    \end{tikzpicture}
  \end{center}

  \begin{block}{Lookup versus learning}
    During a forward pass, lookup simply selects a column. During training,
    gradients update the values in every token column. Over many examples,
    the token vectors develop meaningful geometric relationships based on how
    the tokens are used.
  \end{block}
\end{frame}

\begin{frame}{Training learns the embedding values}
  \small
  The embedding matrix \(W_E\) is a set of model parameters. Token IDs select its
  columns, but next-token training determines the numbers stored in them.

  \vspace{0.4em}
  \begin{center}
    \begin{tikzpicture}[
      stage/.style={draw, rounded corners=3pt, minimum height=1.15cm,
        minimum width=2.35cm, align=center},
      >=Latex
    ]
      \node[stage, fill=tokenblue, draw=tokenblueborder] (initial)
        {\textbf{1. Initialize}\\small random \(W_E\)};
      \node[stage, fill=tokenorange, draw=tokenorangeborder,
        right=0.42cm of initial] (predict)
        {\textbf{2. Predict}\\next token};
      \node[stage, fill=tensorpurple, draw=tensorpurpleborder,
        right=0.42cm of predict] (loss)
        {\textbf{3. Measure}\\loss \(L\)};
      \node[stage, fill=tokengreen, draw=tokengreenborder,
        right=0.42cm of loss] (update)
        {\textbf{4. Update}\\parameters};
      \draw[->, thick] (initial) -- (predict);
      \draw[->, thick] (predict) -- (loss);
      \draw[->, thick] (loss) -- (update);
      \draw[->, thick, rounded corners]
        (update.south) -- ++(0,-0.35) -| (predict.south);
    \end{tikzpicture}
  \end{center}

  \begin{columns}[T]
    \begin{column}{0.47\textwidth}
      \begin{block}{Next-token loss}
        If the correct next token is \(y\), the cross-entropy contribution is
        \[
          L=-\log p(y\mid x).
        \]
        High probability for \(y\) gives a small loss; low probability gives a
        large loss.
      \end{block}
    \end{column}
    \begin{column}{0.47\textwidth}
      \begin{block}{Learning from the loss}
        Backpropagation computes \(\nabla_{W_E}L\). An optimizer with learning rate
        \(\eta{}\) updates the matrix:
        \[
          W_E\leftarrow W_E-\eta\nabla_{W_E}L.
        \]
        Repeating this loop gives the columns useful values.
      \end{block}
    \end{column}
  \end{columns}
\end{frame}

\begin{frame}{Cross-entropy strongly penalizes low probability}
  \small
  Let \(p=p(y\mid x)\) be the probability assigned to the correct next token.
  Its loss is \(L(p)=-\log p\) (using the natural logarithm).

  \begin{center}
    \begin{tikzpicture}[x=9.2cm, y=0.88cm, >=Latex]
      \fill[tokenorange, opacity=0.65] (0,0) rectangle (0.25,3.45);
      \fill[tokengreen, opacity=0.65] (0.75,0) rectangle (1,3.45);

      \draw[->, thick] (0,0) -- (1.08,0)
        node[below right=1pt] {\(p\)};
      \draw[->, thick] (0,0) -- (0,3.55)
        node[above left] {\(L=-\log p\)};

      \foreach \x/\label in {0.1/{0.1},0.5/{0.5},0.9/{0.9},1/{1.0}} {
        \draw (\x,0.06) -- (\x,-0.06) node[below] {\(\label\)};
      }
      \foreach \y in {1,2,3} {
        \draw (0.007,\y) -- (-0.007,\y) node[left] {\(\y\)};
      }

      \draw[very thick, tokenblueborder, domain=0.04:1, samples=100,
        variable=\p]
        plot ({\p},{-ln(\p)});

      \fill[tokenorangeborder] (0.1,2.3026) circle (2.3pt);
      \node[anchor=west, fill=white, inner sep=2pt]
        at (0.13,2.55) {\(p=0.10 \Rightarrow L\approx2.30\)};

      \fill[tokenblueborder] (0.5,0.6931) circle (2.3pt);
      \node[anchor=south west, fill=white, inner sep=2pt]
        at (0.51,0.78) {\(p=0.50 \Rightarrow L\approx0.69\)};

      \fill[tokengreenborder] (0.9,0.1054) circle (2.3pt);
      \node[anchor=south east, fill=white, inner sep=2pt]
        at (0.88,0.18) {\(p=0.90 \Rightarrow L\approx0.11\)};

      \node[tokenorangeborder, anchor=west] at (0.02,3.25)
        {\(p\to0 \Rightarrow L\to\infty\)};
    \end{tikzpicture}
  \end{center}

  \vspace{-0.15em}
  \begin{alertblock}{Training signal}
    A confident wrong prediction receives a large penalty; assigning the
    correct token probability near \(1\) drives its loss toward \(0\).
  \end{alertblock}
\end{frame}

\begin{frame}{Unembedding is needed to make predictions}
  \small
  To evaluate the loss \(L\), the model must make a next-token prediction.
  For now, treat the Transformer as a black box: it turns the input token
  vectors into a context vector \(\mathbf h\).

  \begin{center}
    \begin{tikzpicture}[
      stage/.style={draw, rounded corners=3pt, minimum width=2.65cm,
        minimum height=1.05cm, align=center},
      >=Latex
    ]
      \node[stage, fill=tensorpurple, draw=tensorpurpleborder] (transformer)
        {Transformer\\context \(\mathbf h\)};
      \node[stage, fill=tokenorange, draw=tokenorangeborder,
        right=0.9cm of transformer] (unembed)
        {learned unembedding\\matrix \(W_U\)};
      \node[stage, fill=tokenblue, draw=tokenblueborder,
        right=0.9cm of unembed] (logits)
        {one score per token\\logits \(\mathbf z\)};
      \draw[->, thick] (transformer) -- (unembed);
      \draw[->, thick] (unembed) -- (logits);
    \end{tikzpicture}
  \end{center}

  \begin{columns}[T]
    \begin{column}{0.47\textwidth}
      \begin{block}{Embedding \(W_E\)}
        Token ID \(\longrightarrow{}\) input token vector
      \end{block}
    \end{column}
    \begin{column}{0.47\textwidth}
      \begin{block}{Unembedding \(W_U\)}
        Context vector \(\longrightarrow{}\) token scores
      \end{block}
    \end{column}
  \end{columns}

  \begin{alertblock}{Why training needs both matrices}
    \(W_E\) supplies the vectors that enter the Transformer. \(W_U\) supplies
    one candidate vector for every possible next token, letting the model
    produce the scores needed for its prediction and loss.
  \end{alertblock}
\end{frame}

\begin{frame}{Dot product: scoring a candidate token}
  \small
  Input token vectors come from columns of \(W_E\). After the Transformer has
  combined their context, candidate token \(j\) uses column
  \(\mathbf w_j={(W_U)}_{:,j}\). Its logit is the dot product
  \(z_j=\mathbf{h}^{\mathsf T}\mathbf{w}_j\).

  \begin{columns}[T]
    \begin{column}{0.41\textwidth}
      \centering
      \begin{tikzpicture}[scale=0.72, >=Latex]
        \coordinate (o) at (0,0);
        \draw[->, very thick, tokenblueborder] (o) -- (4.0,0.7)
          node[right] {\(\mathbf{h}\)};
        \draw[->, very thick, tokengreenborder] (o) -- (2.8,2.7)
          node[above] {\(\mathbf{w}_j\)};
        \draw (1.0,0.18) arc[start angle=10,end angle=44,radius=1.05]
          node[midway, right] {\(\theta\)};
      \end{tikzpicture}

      \[
        \mathbf{h}^{\mathsf T}\mathbf{w}_j
        =\lVert\mathbf{h}\rVert
         \lVert\mathbf{w}_j\rVert\cos\theta
      \]
    \end{column}
    \begin{column}{0.55\textwidth}
      \footnotesize
      \begin{block}{Dot-product intuition}
        The score grows when the vectors point in similar directions, but also
        when either vector is longer.
      \end{block}

      \begin{block}{Cosine removes length}
        \[
          \cos\theta=
          \frac{\mathbf{h}^{\mathsf T}\mathbf{w}_j}
          {\lVert\mathbf{h}\rVert\lVert\mathbf{w}_j\rVert}.
        \]
        Dividing out both lengths keeps only directional similarity.
      \end{block}
    \end{column}
  \end{columns}

  \vspace{-0.2em}
  \begin{alertblock}{Connection to cross-entropy}
    \footnotesize
    Dot products produce candidate scores. Cross-entropy pushes the correct
    token's score above its competing tokens.
  \end{alertblock}
\end{frame}

\begin{frame}{From logits to probabilities: softmax}
  \small
  The model produces one raw score, or \emph{logit} \(z_j\), for every token
  \(j\in V\). Logits can be any real numbers and do not have to sum to \(1\).

  \vspace{0.3em}
  \begin{columns}[c]
    \begin{column}{0.49\textwidth}
      \centering
      \textbf{Normalize all vocabulary scores together}
      \[
        p_j=p(y=j\mid x)
        =\frac{e^{z_j}}{\displaystyle\sum_{k\in V}e^{z_k}}
      \]
    \end{column}
    \begin{column}{0.46\textwidth}
      \begin{itemize}
        \item \(e^{z_j}>0\) makes every result positive.
        \item The shared denominator makes
          \(\sum_{j\in V}p_j=1\).
        \item Larger logits receive larger probabilities.
      \end{itemize}
    \end{column}
  \end{columns}

  \begin{center}
    \footnotesize
    \renewcommand{\arraystretch}{1.12}
    \begin{tabular}{cccc}
      \toprule
      Candidate token & Logit \(z_j\) & \(e^{z_j}\) & Probability \(p_j\) \\
      \midrule
      \texttt{cat}  & \(2.0\) & \(7.39\) & \(0.67\) \\
      \texttt{dog}  & \(1.0\) & \(2.72\) & \(0.24\) \\
      \texttt{runs} & \(0.0\) & \(1.00\) & \(0.09\) \\
      \midrule
      & & \(11.11\) & \(1.00\) \\
      \bottomrule
    \end{tabular}
  \end{center}

  \vspace{-1.4em}
  \begin{alertblock}{Connecting probabilities to the loss}
    \footnotesize
    If the correct token is \texttt{cat}, then \(p_y=0.67\) and
    \(L=-\log(0.67)\approx0.40\).
  \end{alertblock}
\end{frame}

\begin{frame}{Why does \(e\) appear in softmax?}
  \small
  \begin{center}
    \large\itshape
    ``That's a great question, because the answer isn't
    \emph{'because mathematicians like \(e\)'}''.
  \end{center}

  \begin{center}
    \begin{tikzpicture}[
      stage/.style={draw, rounded corners=3pt, minimum height=1.0cm,
        minimum width=2.9cm, align=center},
      >=Latex
    ]
      \node[stage, fill=tensorpurple, draw=tensorpurpleborder] (scores)
        {logits\\\(z_j\in\mathbb R\)};
      \node[stage, fill=tokenorange, draw=tokenorangeborder,
        right=0.7cm of scores] (weights)
        {exponentiate\\\(w_j=e^{z_j}>0\)};
      \node[stage, fill=tokengreen, draw=tokengreenborder,
        right=0.7cm of weights] (probabilities)
        {normalize\\\(p_j=\dfrac{w_j}{\sum_k w_k}\)};
      \draw[->, thick] (scores) -- (weights);
      \draw[->, thick] (weights) -- (probabilities);
    \end{tikzpicture}
  \end{center}

  \vspace{-1em}
  \begin{columns}[T,onlytextwidth]
    \begin{column}{0.49\textwidth}
      \begin{block}{Why an exponential?}
        \small
        It is positive, preserves order, and a common shift cancels:
        \[
          \frac{e^{z_j+c}}{\sum_k e^{z_k+c}}
          =\frac{e^c e^{z_j}}{e^c\sum_k e^{z_k}}
          =\frac{e^{z_j}}{\sum_k e^{z_k}}.
        \]
        Only differences between logits affect the probabilities.
      \end{block}
    \end{column}
    \begin{column}{0.45\textwidth}
      \begin{block}{Why specifically base \(e\)?}
        \small
        Any base \(a>1\) could work, since
        \[
          a^z=e^{(\ln a)z}.
        \]
        The base only rescales the logits. But
        \[
          \frac{d}{dz}a^z=(\ln a)a^z,
          \qquad
          \frac{d}{dz}e^z=e^z.
        \]
        Base \(e\) removes the extra factor, keeping learning gradients clean.
      \end{block}
    \end{column}
  \end{columns}
\end{frame}

\begin{frame}{Embeddings turn usage into geometry}
  \small
  \begin{columns}[T]
    \begin{column}{0.55\textwidth}
      \centering
      \begin{tikzpicture}[scale=0.78, >=Latex]
        \draw[->, gray] (0,0) -- (6.0,0);
        \draw[->, gray] (0,0) -- (0,4.5);

        \foreach \point in {
          (1.1,3.4),(2.2,3.5),(1.6,2.7),(2.7,2.8),
          (4.6,1.2),(5.1,1.7)
        } {
          \draw[->, thin, gray!45] (0,0) -- \point;
        }

        \fill[tokenblueborder] (1.1,3.4) circle (2.5pt)
          node[above left]{\texttt{cat}};
        \fill[tokenblueborder] (2.2,3.5) circle (2.5pt)
          node[above]{\texttt{dog}};
        \fill[tokenblueborder] (1.6,2.7) circle (2.5pt)
          node[below left]{\texttt{kitten}};
        \fill[tokenblueborder] (2.7,2.8) circle (2.5pt)
          node[below right]{\texttt{puppy}};

        \draw[->, very thick, tensorpurpleborder] (1.1,3.4) -- (1.6,2.7)
          node[midway, left, font=\scriptsize] {\(\Delta\)};
        \draw[->, very thick, tensorpurpleborder] (2.2,3.5) -- (2.7,2.8)
          node[midway, right, font=\scriptsize] {\(\Delta\)};

        \fill[tokenorangeborder] (4.6,1.2) circle (2.5pt)
          node[below]{\texttt{car}};
        \fill[tokenorangeborder] (5.1,1.7) circle (2.5pt)
          node[above]{\texttt{truck}};

      \end{tikzpicture}

      \vspace{-0.3em}
      {\scriptsize
      \(
        \begin{aligned}
          \mathbf{e}_{\text{puppy}}-\mathbf{e}_{\text{dog}}
            &\approx \boldsymbol{\Delta},\\
          \mathbf{e}_{\text{kitten}}-\mathbf{e}_{\text{cat}}
            &\approx \boldsymbol{\Delta}.
        \end{aligned}
      \)}

      \smallskip
      \textbf{Both arrows show the same relative move:}
      adult animal \(\rightarrow{}\) young animal.
    \end{column}
    \begin{column}{0.42\textwidth}
      Tokens used in similar contexts often acquire related locations.

      \vspace{0.8em}
      Cosine similarity measures directional similarity:
      \[
        \cos(\mathbf{a},\mathbf{b})=
        \frac{\mathbf{a}\cdot\mathbf{b}}
        {\lVert\mathbf{a}\rVert\lVert\mathbf{b}\rVert}.
      \]

    \end{column}
  \end{columns}

  \vspace{-0.3em}
  \begin{alertblock}{Measure, not objective}
    \footnotesize
    Next-token cross-entropy is the training objective; cosine similarity only
    inspects learned geometry. Lookup embeddings are static; contextual
    representations come later.
  \end{alertblock}
\end{frame}

\begin{frame}{Scale example: the GPT-3 175B embedding matrix}
  \small
  The largest GPT-3 model uses a vocabulary of \(50{,}257\) tokens and an
  embedding dimension of \(d=12{,}288\).

  \[
    W_E\in\mathbb R^{12{,}288\times 50{,}257}
    \qquad
    \text{(\(d\) rows, one column per token).}
  \]

  \begin{center}
    \begin{tikzpicture}[>=Latex]
      \fill[tokengreen] (0,0) rectangle (8.3,2.35);
      \draw[tokengreenborder, thick] (0,0) rectangle (8.3,2.35);

      \foreach \x in {0.65,1.3,1.95,6.35,7.0,7.65} {
        \draw[tokengreenborder!55] (\x,0) -- (\x,2.35);
      }
      \foreach \y in {0.47,0.94,1.41,1.88} {
        \draw[tokengreenborder!35] (0,\y) -- (8.3,\y);
      }

      \node[font=\Large] at (4.15,1.38) {\(W_E\)};
      \node at (4.15,0.72) {\(617{,}558{,}016\) learned entries};
      \node at (4.15,2.0) {\(\cdots\)};

      \draw[<->, thick] (0,-0.35) -- (8.3,-0.35)
        node[midway, below] {\(|V|=50{,}257\) token columns};
      \draw[<->, thick] (-0.35,0) -- (-0.35,2.35)
        node[midway, above, rotate=90] {\(d=12{,}288\) rows};
    \end{tikzpicture}
  \end{center}

  \vspace{-0.65em}
  \begin{alertblock}{Parameter count and matching dimensions}
    \footnotesize
    \centering
    \(50{,}257\times12{,}288
      =\boxed{617{,}558{,}016}\approx 618\text{ million parameters}\).

    \smallskip
    For GPT-3, \(W_U\) has the same shape as \(W_E\):
    \(W_U,W_E\in\mathbb R^{12{,}288\times50{,}257}\).
    Equal shape does not mean equal entries.
  \end{alertblock}
\end{frame}

\begin{frame}[plain]
  \centering
  \vfill

  {\Huge Short break}

  \vspace{1.2em}
  {\large Tokens \(\rightarrow{}\) IDs \(\rightarrow{}\) embeddings \(\rightarrow{}\)
  next-token loss \quad \checkmark}

  \vspace{2.0em}
  {\Large Next: inside the Transformer}

  \vspace{0.5em}
  {\normalsize How do token representations exchange context?}

  \vfill
\end{frame}

\begin{frame}{A token sequence becomes a matrix}
  \small
  \begin{center}
    \tokenbox{the}\quad\tokenbox{cat}\quad\tokenbox{sleeps}
    \(\quad\longrightarrow\quad{}\)
    \([12,51,908]\)
  \end{center}

  \vspace{0.3em}
  Each ID selects one column from \(W_E\). Stacking the transposed vectors gives
  \[
    X=
    \begin{bmatrix}
      \mathbf{e}_{12}^{\mathsf T}\\
      \mathbf{e}_{51}^{\mathsf T}\\
      \mathbf{e}_{908}^{\mathsf T}
    \end{bmatrix}
    \in\mathbb{R}^{n\times d},
    \qquad n=3.
  \]

  \begin{center}
    \begin{tikzpicture}[>=Latex, node distance=0.55cm]
      \node[draw=tokenblueborder, fill=tokenblue, rounded corners,
        minimum width=3.1cm] (sequence) {one sequence: \(n\times d\)};
      \node[draw=tensorpurpleborder, fill=tensorpurple, rounded corners,
        right=4.0cm of sequence, minimum width=3.4cm] (batch)
        {batch: \(T\in\mathbb{R}^{b\times n\times d}\)};
      \draw[->, thick] (sequence) -- node[above]{stack \(b\) sequences} (batch);
    \end{tikzpicture}
  \end{center}

  {\scriptsize Sequences in a batch are commonly padded to a shared length;
  a mask tells later layers which positions are padding.}
\end{frame}

\section{Transformers and attention}

\begin{frame}{From token vectors to contextual vectors}
  \begin{center}
    \begin{tikzpicture}[
      stage/.style={draw, rounded corners=3pt, minimum width=2.1cm,
        minimum height=0.8cm, align=center},
      >=Latex
    ]
      \node[stage, fill=tokenblue, draw=tokenblueborder] (tokens) {tokens};
      \node[stage, fill=tokenorange, draw=tokenorangeborder,
        right=0.55cm of tokens] (ids) {IDs};
      \node[stage, fill=tokengreen, draw=tokengreenborder,
        right=0.55cm of ids] (vectors) {embeddings};
      \node[stage, fill=tensorpurple, draw=tensorpurpleborder,
        right=0.55cm of vectors] (transformer) {Transformer\\?};
      \draw[->, thick] (tokens) -- (ids);
      \draw[->, thick] (ids) -- (vectors);
      \draw[->, thick] (vectors) -- (transformer);
    \end{tikzpicture}
  \end{center}

  \vspace{0.8em}
  The embedding lookup gives \tokenbox{bank} the same starting vector in both
  sentences:
  \begin{itemize}
    \item ``They sat on the river \textbf{bank}''.
    \item ``She deposited money at the \textbf{bank}''.
  \end{itemize}

  \begin{alertblock}{The problem self-attention solves}
    How can every token exchange information with the other tokens so that its
    vector reflects the surrounding context? The following slides build the
    mathematical mechanism that makes this possible.
  \end{alertblock}
\end{frame}
